\documentclass{article}

\usepackage[utf8]{inputenc}     % pdfLaTeX
\usepackage[T1]{fontenc}
\usepackage[magyar]{babel}
\usepackage{amsmath}

% Set page size and margins
% Replace `letterpaper' with `a4paper' for UK/EU standard size
\usepackage[a4paper,top=2cm,bottom=2cm,left=3cm,right=3cm,marginparwidth=1.75cm]{geometry}
%for arrows
\usepackage{textcomp}

% --- more levels stuff begins here, dont touch
% --- Section formatting ---
\usepackage{titlesec}

% Allow deep section numbering and TOC
\setcounter{secnumdepth}{5}
\setcounter{tocdepth}{5}

\makeatletter

% ------------------------
% LEVEL 4: subsubsubsection
% ------------------------
\titleclass{\subsubsubsection}{straight}[\subsubsection]
\newcounter{subsubsubsection}[subsubsection]
\renewcommand\thesubsubsubsection{\thesubsubsection.\arabic{subsubsubsection}}

\renewcommand\subsubsubsection{%
  \@startsection{subsubsubsection}{4}{\z@}%
  {3.25ex \@plus1ex \@minus.2ex}%
  {1.5ex \@plus .2ex}%
  {\normalfont\normalsize\bfseries}%
}

\titleformat{\subsubsubsection}
  {\normalfont\normalsize\bfseries}
  {\thesubsubsubsection}{1em}{}

\titlespacing*{\subsubsubsection}{0pt}{3ex}{1ex}

% TOC entry for level 4
\newcommand*\l@subsubsubsection{\@dottedtocline{4}{7em}{4em}}

% --------------------------
% LEVEL 5: subsubsubsubsection
% --------------------------
\titleclass{\subsubsubsubsection}{straight}[\subsubsubsection]
\newcounter{subsubsubsubsection}[subsubsubsection]
\renewcommand\thesubsubsubsubsection{\thesubsubsubsection.\arabic{subsubsubsubsection}}

\renewcommand\subsubsubsubsection{%
  \@startsection{subsubsubsubsection}{5}{\z@}%
  {3.25ex \@plus1ex \@minus.2ex}%
  {1.5ex \@plus .2ex}%
  {\normalfont\normalsize\bfseries}%
}

\titleformat{\subsubsubsubsection}
  {\normalfont\normalsize\bfseries}
  {\thesubsubsubsubsection}{1em}{}

\titlespacing*{\subsubsubsubsection}{0pt}{3ex}{1ex}

% TOC entry for level 5
% Use level 4 internally for safety (LaTeX only knows TOC levels 1-4)
\newcommand*\l@subsubsubsubsection{\@dottedtocline{4}{10em}{5em}}

% --------------------------
% LEVEL 6: subsubsubsubsubsection
% --------------------------
\titleclass{\subsubsubsubsubsection}{straight}[\subsubsubsubsection]

\newcounter{subsubsubsubsubsection}[subsubsubsubsection]
\renewcommand\thesubsubsubsubsubsection{%
  \thesubsubsubsubsection.\arabic{subsubsubsubsubsection}}

\renewcommand\subsubsubsubsubsection{%
  \@startsection{subsubsubsubsubsection}{6}{\z@}%
  {3.25ex \@plus1ex \@minus.2ex}%
  {1.5ex \@plus .2ex}%
  {\normalfont\normalsize\bfseries}%
}

\titleformat{\subsubsubsubsubsection}
  {\normalfont\normalsize\bfseries}
  {\thesubsubsubsubsubsection}{1em}{}

\titlespacing*{\subsubsubsubsubsection}{0pt}{3ex}{1ex}

% TOC entry (still forced into level 4)
\newcommand*\l@subsubsubsubsubsection{\@dottedtocline{4}{13em}{6em}}

% --- more levels stuff ends here

\title{Mássalhangzótörvények}
\author{Lengyel Zoltán Péter\thanks{Csukás Ibolya tanítása alapján}}
\date{Legutóbb frissítve: 2026. Május 05.}

\begin{document}

\maketitle
\tableofcontents
\newpage

\begin{itemize}
    \item Egymás mellé kerülő mássalhangzók gyakran hatnak egymás képzésére: módosítják vagy megváltoztatják azt; többnyire kiejtésben, ritkábban írásban is
\end{itemize}

\section{Hasonulás}
\begin{itemize}
    \item Két szomszédos mássalhangzó közül az egyik részben vagy teljesen hasonul a másikhoz
\end{itemize}

\subsection{Részleges hasonulás}
\subsubsection{Zöngésség szerinti részleges hasonulás}

\begin{itemize}
    \item zöngés + zöngétlen = zöngétlen  \begin{itemize}
     \item zse\textbf{b}\textit{k}endő $\rightarrow$ [zse\textbf{p}\textit{k}endő] \begin{itemize}
         \item \textbf{b: zöngés} 
         \item \textit{k: zöngétlen}
         \item zöngés + zöngétlen = zöngétlen
         \item \textbf{b $\rightarrow$ p}
     \end{itemize}
 \end{itemize}

    \item zöngétlen + zöngés = zöngés
     \begin{itemize}
     \item mo\textbf{s}\textit{d}ó $\rightarrow$ [mo\textbf{zs}\textit{d}ó] \begin{itemize}
         \item \textbf{s: zöngétlen} 
         \item \textit{d: zöngés}
         \item zöngétlen + zöngés = zöngés
         \item \textbf{s $\rightarrow$ zs}
     \end{itemize}
 \end{itemize}
\end{itemize}

\subsubsection{Képzés helye szerinti részleges hasonulás}

\begin{itemize}
    \item nb $\rightarrow$ mb\begin{itemize}
        \item azo\textbf{n}\textit{b}an $\rightarrow$ [azo\textbf{m}\textit{b}an] \begin{itemize}
            \item \textbf{n: foghang}
            \item \textit{b: ajakhang}
            \item \textbf{n $\rightarrow$ m (ajakhang)}
        \end{itemize}
    \end{itemize}
    \item np $\rightarrow$ mp \begin{itemize}
        \item szí\textbf{n}\textit{p}ad $\rightarrow$ [szí\textbf{m}\textit{p}ad] \begin{itemize}
            \item \textbf{n: foghang}
            \item \textit{p: ajakhang}
            \item \textbf{n $\rightarrow$ m (ajakhang)}
        \end{itemize}
    \end{itemize}
    \item ngy $\rightarrow$ nygy \begin{itemize}
        \item ro\textbf{n}\textit{gy}os $\rightarrow$ [ro\textbf{ny}\textit{gy}os] \begin{itemize}
            \item \textbf{n: foghang}
            \item \textit{gy: szájpadláshang}
            \item \textbf{n $\rightarrow$ ny (szájpadláshang)}
        \end{itemize}
    \end{itemize}
    \item nty $\rightarrow$ nyty \begin{itemize}
        \item po\textbf{n}\textit{ty} $\rightarrow$ [po\textbf{ny}\textit{ty}] \begin{itemize}
            \item \textbf{n: foghang}
            \item \textit{ty: szájpadláshang}
            \item \textbf{n $\rightarrow$ ny (szájpadláshang)}
        \end{itemize}
    \end{itemize}
\end{itemize}




\subsection{Teljes hasonulás}


\subsubsection{Írásban jelöletlen teljes hasonulás}
\begin{itemize}
    \item Két szomszédos mássalhangzó közül az egyik teljesen hasonul a másikhoz, de csak kiejtésben! \begin{itemize}
        \item pl. teljes $\rightarrow$ [tejjes]
        \item adta $\rightarrow$ [atta]
        \item tanmenet $\rightarrow$ [tammenet]
        \item egészség $\rightarrow$ [egésség]
        \item évforduló $\rightarrow$ [évforduló]
        \item fogkrém $\rightarrow$ [fokkrém]
        \item mesélj $\rightarrow$ [meséjj]
    \end{itemize}
\end{itemize}


\subsubsection{Írásban jelölt teljes hasonulás}
\begin{itemize}
    \item Két szomszédos mássalhangzó közül az egyik teljesen hasonul a másikhoz, kiejtésben és írásban is
\end{itemize}

\subsubsubsection{-val/-vel, -vá/vé}
\begin{itemize}
    \item Mássalhagzóra végződő szóhoz -val/-vel/-vá/-vé rag járul; a rag v hangja teljesen hasonul a szó utolsó mássalhangzójához
    \begin{itemize}
        \item Gábor + -val = Gáborral
        \item sajt + -vá = satjjá
    \end{itemize}
\end{itemize}
\subsubsubsubsection{Két különböző mássalhangzóra vegződő szavak} \begin{itemize}
        \item Ha két különböző mássalhangzóra végződik a szó; a hasonult mássalhangzót röviden ejtjük \begin{itemize}
        \item Zsolttal $\rightarrow$ [Zsoltal]
    \end{itemize}
\end{itemize}
\subsubsubsubsection{Hosszú mássalhangzóra végződő szavak} \begin{itemize}
    \item Ha hosszú mássalhangzóra végződik a szó  egyszerűsítünk \begin{itemize}
        \item sakk + -val = sakkal 
    \end{itemize}
    \item DE! Hosszú mássalhangzóra végződő családnevekben és tulajdonnevekben nem egyszerűsíthetünk! Az ellentmondást kötőjellel oldjuk \begin{itemize}
        \item Mariann + -nal = Mariann-nal
        \item Papp + -val = Papp-pal
    \end{itemize}
\end{itemize}


\subsubsubsubsection{Régies betűre végződő nevek}  \begin{itemize}
        \item Régies betűre végződő nevekben kiejtés szerint kapcsoljuk a ragot
    \begin{itemize}
        \item Kossuth + -val = Kossuthtal
        \item Damjanich + -val = Damjaichcsal
    \end{itemize}
\end{itemize}

\subsubsubsubsection{Néma hangra végződő nevek}
\begin{itemize}
    \item Néma hangra végződő nevekhez kötőjellel kapcsoljuk a kiejtés szerint hasonult ragot \begin{itemize}
        \item Joe + val = Joe-val
        \item Anne + -vel = Anne-nel
        \item Shakespeare + -vel = Shakespeare-rel
    \end{itemize}
\end{itemize}

\subsubsubsubsection{Számok, rövidítések}
\begin{itemize}
    \item Számokhoz, rövidítésekhez kötőjellel kapcsoljuk a kiejtés szerint hasonult ragot \begin{itemize}
        \item 3 + -val = 3-mal
        \item km + -vel = km-rel
        \item \% + -val = \%-kal
    \end{itemize}
\end{itemize}

\subsubsubsubsection{h-ra végződő szavak}
\begin{itemize}
    \item Néhány h-ra végződő szót kétféleképpen is toldalékolhatunk \begin{itemize}
        \item juh + -val = juhhal/juhval
        \item cseh + -vel = csehhel/csehvel
        \item még + -vel = méhhel/mével
    \end{itemize}
\end{itemize}

\subsubsubsection{s, sz, z, dz végű ige + j}
\begin{itemize}
    \item s, sz, z, dz végű igéhez járul: \begin{itemize}
    \item a felszólító mód jele -j \begin{itemize}
        \item mos + -j = moss!
        \item húz + -j = húzz!
        \item edz + -j = eddz!
        \item játsz + -j = játssz!
    \end{itemize}
    \item -j-vel kezdődő személyrag \begin{itemize}
        \item mos + -ja = mossa
        \item játsz + -juk = játsszuk
        \item húz + -játok = húzzátok
        \item edz + -jük = eddzük
    \end{itemize}
    \end{itemize}
\end{itemize}

\subsubsubsection{ez/az mutató névmás + mássalhangzóval kezdődő toldalék}
\begin{itemize}
    \item Ha ez/az  utató névmáshoz mássalhangzóval kezdődő toldalék járul; a z teljesen hasonul a toldalék első mássalhangzójához
    \begin{itemize}
        \item ez + -ben = ebben
        \item az + -tól = attól
    \end{itemize}
\end{itemize}




\subsection{Összeolvadás}
\begin{itemize}
    \item Két szomszédos mássalhangzó helyezz egy harmadik mássalhangzót ejtünk
    \item leggyakrabban összeolvadó mássalhangzók: \begin{itemize}
        \item ts [cs] pl. tatíts, gyógyítson
        \item ds [cs] pl. szabadság
        \item tj [ty] pl. gyógyítja, takarítjuk
        \item dj [gy] pl. maradj, adjuk
        \item ny [ny] pl. fonja, kenjétek
        \item tsz [c] pl. metsz, ötször
        \item dsz [c] pl. adsz, maradsz
    \end{itemize}
\end{itemize}

\subsection{Mássalhangzó-rövidülés}
\begin{itemize}
    \item Ha egy hosszú mássalhangzó előtt vagy után egy másik mássalhangzó áll, a hosszú mássalhangzót röviden ejtjük \begin{itemize}
        \item Zso\underline{l}\textbf{tt}al
        \item vo\underline{n}\textbf{zz}a
        \item o\textbf{tt}\underline{h}on
    \end{itemize}
\end{itemize}

\subsection{Mássalhangzó-kiesés}
\begin{itemize}
    \item Három szomszédos mássalhangzó közül a középső d kiejtésben kiesik \begin{itemize}
        \item mon\underline{d}tam $\rightarrow$ [montam]
        \item kül\underline{d}tétek $\rightarrow$ [kültétek]
    \end{itemize}
\end{itemize}

\end{document}